% MT011 / explanation — canonical editable content.
过拟合是模型把训练样本中的偶然变化也当成了普遍规律。一个极深的树可以不断切分，直到几乎给每个训练样本单独安排一个区域；它容易把训练题答对，但这些很小的区域未必适合新样本。


训练集用于学习模型参数；验证集检查模型对未参与这次拟合的数据表现如何。分类错误率就是分错的条数除以总条数。判断过拟合时，比较同一指标下的训练误差与验证误差。例如训练错1\%、验证错20\%，值得检查过拟合；训练和验证都错20\%，则更像还没学到规律。比较之前要确认两批数据的处理方式和来源大体一致，避免把分布变化误当成纯粹的过拟合。


交叉验证（cross-validation，CV）在训练数据内部轮流留出一部分做验证，用这些验证成绩比较不同方案；五折示例见\StudyRef{Answers}{mt004}{MT004}。可以尝试以下改进：



\begin{enumerate}[start=1]

\item \textbf{限制模型复杂度。} 减小树深、降低多项式次数，或用更简单的模型，让模型不能为每个训练点单独绕路。

\item \textbf{加强合适的正则化。} 给大权重增加代价，并用训练集内部的CV选强度；惩罚过强也会把真实规律压掉。

\item \textbf{减少无关特征。} 无关输入给模型更多凑巧拟合噪声的机会，特征筛选也须放在训练流程内完成。

\item \textbf{增加有代表性的训练数据。} 让同一个真实规律在更多样本中反复出现，比重复收集同一小群体的数据更有帮助。

\end{enumerate}


最终测试集留到方法选择完成后使用。


\textbf{出处：} 2020B期中Quiz Q11，原卷第5页（10分）


\textbf{答案依据：} 2020B期中 Q11 随卷／配套参考答案；条件与解释经复核。
